\documentclass[10pt,letterpaper]{article}
\usepackage{cornell-notes}

\title{Clause 25.05.05.03: Constructors And Conversions}
\author{}
\date{}
\setDocTitle{Clause 25.05.05.03: Constructors And Conversions}
\setDocSubtitle{C++ 2024}
\setDocOwner{C++ 2024 Cornell Notes}
\setDocDate{\today}
\setCornellCollection{C++ 2024 Cornell Notes}
\setCornellUnitType{Clause}
\setCornellUnitNumber{25.05.05.03}
\setCornellUnitTitle{Constructors And Conversions}
\hypersetup{pdftitle={Clause 25.05.05.03: Constructors And Conversions},pdfsubject={Cornell notes},pdfauthor={}}

\begin{document}
\maketitle
\begin{CornellOverviewBox}{Overview}
\textbf{Classification:} Standard library - Normative\\[2pt]
\textbf{Learning objective:} Understand the rules, terminology, and semantic consequences associated with Constructors and conversions.
\end{CornellOverviewBox}\begin{CornellSummaryBox}{Summary}
{\sffamily\bfseries\color{Accent} Summary}\par\vspace{3pt}Section 25.5.5.3 focuses on Constructors and conversions. Apply the specified interface together with its mandates, effects, result conditions, exception behavior, and complexity constraints. When applying it, preserve the distinction between mandatory requirements, recommendations, permissions, and behavior deliberately left undefined, unspecified, or implementation-defined.
\end{CornellSummaryBox}

\end{document}
